% 第 6 章 全球竞争格局与国产替代进度
% 三大原厂 HBM 份额饼 + 国产替代进度总表

\section{三大原厂 HBM 格局：SK 海力士绝对领先}

HBM 市场高度集中，呈现「一强两弱」格局。SK 海力士凭借先发优势 + 与 NVIDIA 的深度技术合作，自 HBM2E 时代起持续保持龙头地位。HBM3E 的良率差距 (SK 85-88\% vs 三星 78-82\% vs 美光 75-80\%) 决定了 2026-2027 年的份额分配。

\needspace{44\baselineskip}
\begin{exhibitbox}[图 6-1 \quad 三大原厂 HBM 产能份额预测 (2025E / 2026E / 2027E) / HBM Capacity Share]
\centering
\resizebox{0.88\exhibitwidth}{!}{%
\begin{tikzpicture}[
  pie/.style={circle, minimum size=4.0cm},
  slice/.style={insert path={-- ++(\a:2.0) arc (\a:\a+#1:2.0) -- cycle}},
  slabel/.style={font=\bfseries\small, align=center, text=deepnavy}
]
% Pie 1: 2025E
\begingroup
\node[above=0.1cm, font=\bfseries\large, color=navy] at (0, 2.3) {\textbf{2025E}};
% SK 60% = 216deg, Samsung 24% = 86.4deg, Micron 12% = 43.2deg, Other 4% = 14.4deg
\def\a{90}
\filldraw[fill=riskred!80, draw=white, line width=1.2pt] (0,0) [slice={216}];
\pgfmathsetmacro\a{\a+216}
\filldraw[fill=accentblue!80, draw=white, line width=1.2pt] (0,0) [slice={86.4}];
\pgfmathsetmacro\a{\a+86.4}
\filldraw[fill=nvgreen!75, draw=white, line width=1.2pt] (0,0) [slice={43.2}];
\pgfmathsetmacro\a{\a+43.2}
\filldraw[fill=steelgrey!60, draw=white, line width=1.2pt] (0,0) [slice={14.4}];
\node[slabel, text=white] at (30:1.3) {SK 60\%};
\node[slabel, text=white] at (240:1.2) {Samsung 24\%};
\node[slabel, text=white] at (305:1.2) {Micron 12\%};
\node[slabel, text=white, font=\tiny] at (345:1.5) {4\%};
\draw (0,0) circle (2.0);
\endgroup

% Pie 2: 2026E
\begin{scope}[xshift=6.0cm]
\begingroup
\node[above=0.1cm, font=\bfseries\large, color=navy] at (0, 2.3) {\textbf{2026E}};
\def\a{90}
\filldraw[fill=riskred!80, draw=white, line width=1.2pt] (0,0) [slice={205.2}];
\pgfmathsetmacro\a{\a+205.2}
\filldraw[fill=accentblue!80, draw=white, line width=1.2pt] (0,0) [slice={97.2}];
\pgfmathsetmacro\a{\a+97.2}
\filldraw[fill=nvgreen!75, draw=white, line width=1.2pt] (0,0) [slice={54.0}];
\pgfmathsetmacro\a{\a+54.0}
\filldraw[fill=steelgrey!60, draw=white, line width=1.2pt] (0,0) [slice={3.6}];
\node[slabel, text=white] at (33:1.3) {SK 57\%};
\node[slabel, text=white] at (243:1.2) {Samsung 27\%};
\node[slabel, text=white] at (315:1.2) {Micron 15\%};
\node[slabel, text=white, font=\tiny] at (355:1.7) {1\%};
\draw (0,0) circle (2.0);
\endgroup
\end{scope}

% Pie 3: 2027E
\begin{scope}[xshift=12.0cm]
\begingroup
\node[above=0.1cm, font=\bfseries\large, color=navy] at (0, 2.3) {\textbf{2027E}};
\def\a{90}
\filldraw[fill=riskred!80, draw=white, line width=1.2pt] (0,0) [slice={187.2}];
\pgfmathsetmacro\a{\a+187.2}
\filldraw[fill=accentblue!80, draw=white, line width=1.2pt] (0,0) [slice={108.0}];
\pgfmathsetmacro\a{\a+108.0}
\filldraw[fill=nvgreen!75, draw=white, line width=1.2pt] (0,0) [slice={64.8}];
\node[slabel, text=white] at (40:1.3) {SK 52\%};
\node[slabel, text=white] at (238:1.2) {Samsung 30\%};
\node[slabel, text=white] at (320:1.2) {Micron 18\%};
\draw (0,0) circle (2.0);
\endgroup
\end{scope}

% Legend (bottom)
\node[draw=rulegrey, fill=white, rounded corners=3pt, font=\small, inner sep=5pt, anchor=north] at (6.0, -3.0) {
\begin{tabular}{c l c l c l}
\cellcolor{riskred!60}\rule{0pt}{8pt}\rule{10pt}{0pt} & SK Hynix &
\cellcolor{accentblue!60}\rule{0pt}{8pt}\rule{10pt}{0pt} & Samsung &
\cellcolor{nvgreen!55}\rule{0pt}{8pt}\rule{10pt}{0pt} & Micron \\
\cellcolor{steelgrey!60}\rule{0pt}{8pt}\rule{10pt}{0pt} & Other（铠侠/Solidigm）& \\[2pt]
\end{tabular}
};
\end{tikzpicture}
}
\end{exhibitbox}
\sourcenote{份额数据：SK Hynix 2026Q1 CC [S-11] + DigiTimes 产能报道 + {\faExclamationTriangle} AStock 基于 capex 结构/良率推算。三饼合计 100\%，忽略未上市铠侠/Solidigm 等小份额厂商。}

\vspace{0.4em}
\noindent\textbf{投资含义}：三大原厂 HBM 格局的高度集中意味着 SK 海力士在 2026-2027 年仍将是 HBM 供应链的定价者。A 股在 HBM 原厂层的缺失反而强化了一个核心投资逻辑——\textbf{无论 SK、三星还是美光谁的份额变化，三家原厂都需要扩张 HBM 产线，因此半导体设备 (北方华创/中微/拓荆) 始终是全链条确定性最高的受益者}。先进封装层面，TSMC 的 CoWoS 主导地位短期不可撼动，长电/通富需等待 TSMC 产能满载后的订单外溢 (预计 2027H2+)。

\section{国产替代阶梯：NOR→DRAM→NAND→HBM，HBM 仍差 3-4 代}

国产存储替代遵循清晰的技术难度阶梯：NOR Flash 已实现全球竞争力 (兆易全球前三)；利基 DRAM (长鑫 17nm) 已在消费级量产；3D NAND (长存 232L) 已进入消费级主流；\textbf{HBM 层级仍处于预研阶段，实质可用最早 2028 年后}。

\begin{exhibitbox}[表 6-1 \quad 国产替代进度总表 / Domestic Substitution Progress Matrix]
\centering
\scriptsize
\renewcommand{\arraystretch}{1.2}
\begin{tabularx}{\exhibitboxwidth}{>{\bfseries}p{2.1cm} >{\raggedright\arraybackslash\hspace{0pt}}p{1.9cm} >{\centering\arraybackslash}p{1.8cm} >{\centering\arraybackslash}p{1.2cm} >{\centering\arraybackslash\hspace{0pt}}p{1.3cm} >{\raggedright\arraybackslash\hspace{0pt}}X}
\toprule
\textbf{国内厂商} & \textbf{最新制程节点} & \textbf{月产能 (K 片 12 寸)} & \textbf{良率水平} & \textbf{BIS 约束} & \textbf{客户导入 + A 股关联} \\
\midrule
\rowcolor{nvgreen!8}
\textbf{长江存储 (YMTC)} & 232L TLC/QLC 量产\par \textcolor{steelgrey}{290L 研发中 (2026 试产)} & 280--320K & 232L TLC 85--90\%\par QLC 75--80\% & \textbf{高}\par DUV+刻蚀受限 & 消费级国内渗透 约15\%；企业级小批量 + 华为深度导入。\par 未上市；供应商：\textbf{北方华创/中微/拓荆/沪硅} ★ ★ ★ \\
\midrule
\rowcolor{nvgreen!8}
\textbf{长鑫存储 (CXMT)} & 17nm (1α-) 量产\par \textcolor{steelgrey}{15nm (1β-) 2026 试产} & 220--260K & 17nm DDR4 88--92\%\par DDR5\slash LPDDR5 75--80\% & \textbf{中高}\par 15nm 以下受限 & 消费级 DDR4 国内 约20\%；小米\slash OPPO\slash 传音导入；服务器 DDR5 认证中。\par 未上市；供应商：\textbf{北方华创/深科技/兆易 (潜在 IP 合作)} \\
\midrule
\rowcolor{accentblue!8}
\textbf{兆易创新 (603986)} & NOR: 38nm 量产，2xnm 研发\par \textcolor{steelgrey}{DRAM: 19nm DDR5 试产 + 17nm 研发} & NOR 120--150K (8+12 寸)\par DRAM {\faExclamationTriangle}<5K & NOR 92--95\%\par DRAM {\faExclamationTriangle}<40\% & \textbf{中低}\par 当前节点暂未受限 & \textbf{NOR 全球前三 (约15\%)}，汽车电子/工业客户广；自研 DRAM 消费级小批量送样。\par \textbf{直接上市标的}。 \\
\midrule
\textbf{北京君正 (300223)} & DRAM (ISSI) 38--25nm 利基\par SRAM 全球前三 & 40--50K (代工 Tower\slash UMC) & 成熟制程 >90\% & \textbf{低}\par 利基市场非管制重点 & 汽车\slash 工业\slash 医疗利基 DRAM\allowbreak+\allowbreak SRAM 稳定出货；AI 敞口实质小 (\(<5\%\))。\par 直接上市标的。 \\
\midrule
\rowcolor{riskamber!8}
\textbf{\faExclamationTriangle 国产 HBM 综合评估} & \textcolor{riskred}{\textbf{预研阶段，无公开样品}} & N/A & N/A & \textbf{极高}\par HBM 管制传闻 & 长存/长鑫未披露任何 HBM 计划；2028 年前实质商用概率 <30\%。\par \textbf{A 股无直接受益标的}。 \\
\midrule
\rowcolor{panelgrey}
\multicolumn{2}{>{\cellcolor{panelgrey}\bfseries\raggedright\arraybackslash}p{\dimexpr2.1cm+1.9cm+2\tabcolsep\relax}}{\textbf{设备材料国产替代 (整体)}} & \multicolumn{2}{>{\cellcolor{panelgrey}\centering\arraybackslash}p{\dimexpr1.8cm+1.2cm+2\tabcolsep\relax}}{长存/长鑫 2026E 国产率 约20--30\% (2023 年仅 10\%)} & \cellcolor{panelgrey}\centering\arraybackslash\textbf{核心矛盾} & \cellcolor{panelgrey}\raggedright\arraybackslash\hspace{0pt} BIS 管制是国产化率的双重驱动：短期抑制原厂扩产节奏，长期加速国产设备材料验证周期。\\
\bottomrule
\end{tabularx}
\end{exhibitbox}
\sourcenote{YMTC/CXMT：集微网+地方政府信息 [S-09]；兆易：2025 年报 + IR [S-16]；设备国产率：招标网公开数据 [S-10]。{\faExclamationTriangle}：AStock 行业推断，非原厂官方披露。}

\vspace{0.4em}
\noindent\textbf{投资含义}：国产替代进度总表揭示三个关键定价差异。\textbf{(1) 设备/材料层：国产化率从 10\% \(\to\) 25-30\% 的「S 曲线跃升期」}，北方华创、中微、拓荆、沪硅的收入增速将显著高于原厂 capex 增速 (30-40\% vs 10-30\%)。\textbf{(2) 原厂设计层：A 股在 HBM 上存在重大预期错配}——市场往往将长存/长鑫的制程进展线性外推到 HBM，但 HBM 需要 CoWoS 先进封装 (TSV 工艺) 与原厂的协同设计能力，国内这两项能力均落后 3 代以上。\textbf{(3) 利基存储层：兆易创新/北京君正处于国产替代的「收获期」}，利基 DRAM 因海外大厂主动退出呈现结构性供给收缩，涨价确定性高于标准品。

\section{BIS 管制：国产替代的「双刃剑」}

BIS 出口管制对国产存储链的影响是高度非对称的：对原厂层 (YMTC/CXMT) 短期构成供给约束 (产能扩张受限 12-18 个月)，但对设备/材料/接口芯片层形成长期的「强制国产化」驱动。传闻中的 2026Q3 HBM 单独禁售将进一步放大 CXL 池化内存在国内的战略价值。

从历史管制路径看，BIS 自 2022 年 10 月首次将 YMTC 列入 Entity List 以来，管制呈现「阶梯式收紧、留有灰色空间」的特征：2024 年 10 月修订将 18nm 以下 DRAM 设备纳入禁售范围，但 KrF 光刻和部分成熟刻蚀仍留有余地；2025 年 10 月的进一步修订则将 DUV 光刻机的对华许可审批从「一般审查」升级为「推定拒绝」，直接影响了长存武汉二期和长鑫北京新厂的设备交付节奏。2026 年 5 月路透社报道的「HBM3E/HBM4 单独禁售」传闻若在 Q3 落地，将对国内 AI 芯片厂 (华为昇腾、寒武纪、海光) 的高端产品线形成直接冲击，但同时也会加速国内 CXL 池化内存和企业级 SSD 的替代路径落地。

对 A 股投资而言，BIS 管制本质上是「短期利空原厂 (产能扩张)、长期利多设备材料 (国产化率强制提升)」的双刃剑结构。配置含义有三：第一，若 Q3 管制落地且超预期 (DUV 全面断供)，应短期规避长存/长鑫供应链敞口集中的标的 (部分设备厂订单延后 1-2 季)，但利用回调加仓国产替代逻辑最硬的北方华创 (刻蚀/薄膜平台化) 和中微 (介质刻蚀全球竞争力)；第二，若管制未落地或力度低于预期，则长存/长鑫的扩产加速将直接利好拓荆 (ALD)、盛美上海 (清洗/电镀) 和安集科技 (CMP) 等弹性设备材料股；第三，无论管制是否升级，澜起科技 (CXL 对冲工具) 和江波龙/深科技 (企业级模组国产替代) 都处于管制受益的结构路径上。\par
\textbf{核心监控指标}：BIS Federal Register 每日更新 (HBM 单独管制条目)、ASML 季度订单日志 (对华 DUV 交付数量)、长存设备中标公告月度频次 (中国国际招标网)、华为昇腾 910C 出货量传闻。以上四项每季度交叉验证一次，作为 BIS 管制情景切换的触发信号。
